% problem-000315 5 荧光探针检测一氧化氮
\section*{5. 荧光探针检测一氧化氮}
最近，化学家合成了⼀组荧光分⼦ $\mathrm { ( F L _ { 1 } \sim F L _ { 5 } ) }$ ，可⽤于⽣物体内NO的实时检测。先将荧光分⼦与 $\mathrm { C u C l _ { 2 } }$ 制成配合物，当体系中有NO出现时， $\mathrm { N O }$ 与配合物相互作⽤，从中置换出铜(I)，通过对⽐反应前后体系荧光光谱的变化，即可检测 $\mathrm { N O } _ { \circ }$ 原理⽰意如下（FL代表荧光配体）：

$
\mathrm{FL} + \mathrm{Cu(II)} \rightleftharpoons [ \mathrm{Cu(II)FL} ] \xrightarrow {\mathrm{NO}} \mathrm{Cu(I)} + \mathrm{FL-NO}
$

合成的荧光配体 $\mathrm { F L _ { 1 } }$ ⻅右图：

\subsection*{5-1}
已知 $\mathrm { F L _ { 1 } }$ 与CuCl_{2}⽣成的配合物分⼦式为 $\mathrm { C _ { 3 0 } H _ { 1 8 } N _ { 2 } O _ { 5 } C l _ { 2 } C u }$ ，其中铜为4配位，试写出该配合物的结构式。

\subsection*{5-2}
NO与5-1中的配合物作⽤，⽣成分⼦式为 $\mathrm { C _ { 3 0 } H _ { 1 8 } N _ { 3 } O _ { 6 } C l }$ 的化合物和$\mathrm { C u C l }$ ，试写出该化合物最可能的结构式。

\subsection*{5-3}
荧光分⼦ $\mathrm { F L } _ { 5 } \left( \mathrm { C } _ { 3 2 } \mathrm { H } _ { 2 1 } \mathrm { N } _ { 2 } \mathrm { O } _ { 7 } \mathrm { C l } \right)$ 的合成路线如下，请写出A、B、C、

（图：）

D、E以及 $\mathrm { F L } _ { 5 }$ 的结构式。

()

（图：）

()

（图：）
